%! Piecewise & Step Functions — Unit 1, Lesson 1.4 — Guided Notes & Practice
% THE in-class document, in the MAIN IDEAS / NOTES shape (LESSON_SHAPE.md §1, ported from AP
% Statistics 2026-09-12): the vocabulary box, the hook, then ONE two-column guidednotes table.
% Each row is one idea — a label on the left; on the right one or two COMPLETE printed
% sentences, ONE large pre-drawn display, then a two-across grid of numbered problems with real
% writing room. Problems run 1..13 through the component. No group activity: the notes carry
% the whole gradual release, closed by a SPOKEN debrief; the homework is the individual practice.
%   I DO   : the teacher reads each row's definition, marks up its display, and works the FIRST
%            problem of each grid (1, 3, 5, 7).
%   WE DO  : the class works the rest of each grid, students holding the pen; the last row,
%            Guided Practice (10–13), is worked entirely together on a rule with no story.
%   YOU DO : nothing on this page — ../homework/ IS the individual practice, started in the
%            last ten minutes of the period.
% DENSITY: no mid-sentence blanks — the only \blank{}s are the four cells of the round-down
% table (9) and its one-word follow-up. Every display is pre-drawn and read or labelled
% (\labelbox); no sketch questions.
% ONE CONTEXT: a streaming service, free for the first 3 months, then $12 a month. Every problem
% is "which rule does this month use, and what does the graph do at month 3?"
% THE TRAP is 9's last cell, ⌊-2.5⌋ = -3 (down means more negative). THE CRUX is 7 — exactly
% one piece owns the boundary: P(3) = 0, the closed dot says so; the story cannot decide it, the
% inequality does — and 12 tests it again on new ground.
% PAGE LOCKSTEP with notes_key/: \vterm -> \vtermans, \blank -> \ans, and the answer argument of
% every \pcell and \labelbox filled; every \begin{work} block (the two halves of the trail H) is
% BYTE-IDENTICAL in both files. Nothing else differs. TABLE MECHANICS: inside a cell `\\` ENDS
% THE ROW — break lines with \par; the figure and EACH grid row sit in their own sub-row (`\\`
% then `& ...`, the grid reopening as probgrid*) so the table can break between them.
\documentclass[12pt]{article}
\usepackage{algebra2-article}
\usepackage{algebra2-boxes}
% Vocabulary rows (SAAR style, user correction 2026-09-06): a FIXED-HEIGHT row carrying the term
% label and OPEN writing space — no inline blank and no rule line; students write beside and
% below the term. The key fills exactly the same height (definitions two lines at most), so the
% box cannot drift blank vs. keyed. algebra2-article's \termblank draws a rule line, so the pair
% is defined here; \termrowheight/\termrowinset are taken by the package, hence the local name.
\newcommand{\vrowheight}{2.0cm}
\newcommand{\vterm}[1]{%
  \par\noindent\begin{minipage}[t][\vrowheight][t]{\linewidth}%
    \vspace*{6pt}\textbf{\textcolor{forest}{#1:}}%
  \end{minipage}\par}
% Coordinate grid: {xmin}{xmax}{ymin}{ymax}
\newcommand{\lgrid}[4]{%
  \draw[step=1, linegray!40, very thin] (#1,#3) grid (#2,#4);
  \draw[->, semithick, charcoal] (#1,0)--(#2,0) node[below right, font=\scriptsize]{$x$};
  \draw[->, semithick, charcoal] (0,#3)--(0,#4) node[left, font=\scriptsize]{$y$};}
% Months grid for the streaming service: x = months 0..7, one y-unit = $12.
\newcommand{\mgrid}{%
  \draw[step=1, linegray!40, very thin] (0,0) grid (7,4);
  \draw[->, semithick, charcoal] (0,0)--(7.6,0) node[below right, font=\scriptsize]{$m$};
  \draw[->, semithick, charcoal] (0,0)--(0,4.6) node[left, font=\scriptsize]{\$};
  \foreach \x in {1,...,7}{\node[below, font=\tiny, charcoal] at (\x,0){$\x$};}
  \foreach \y/\lab in {1/12,2/24,3/36,4/48}{\node[left, font=\tiny, charcoal] at (0,\y){$\lab$};}}
% Endpoint dots: {color}{x}{y}. Closed = filled; open = white-filled ring.
\newcommand{\fdot}[3]{\fill[#1] (#2,#3) circle (3.2pt);}
\newcommand{\hdot}[3]{\draw[very thick, #1, fill=white] (#2,#3) circle (3.2pt);}

\begin{document}

\pageheader{Unit 1, Lesson 1.4}{Guided Notes \& Practice}
% NAMESTRIP: no \namedateperiod here — the name row lives on the cover only.

\vspace{6pt}

% ── PAGE 1: vocabulary, then the hook ────────────────────────────────────────
\begin{vocabbox}
\small
Fill in each term \textbf{as we name it} in the notes below.
\vterm{Piecewise-defined function}
\vterm{Boundary (breakpoint)}
\vterm{Open vs.\ closed endpoint}
\vterm{Continuous at a boundary / discontinuity (jump)}
\vterm{Greatest-integer (step) function $\lfloor x\rfloor$}
\end{vocabbox}

\vspace{8pt}

\boxguard[12]
\begin{hookbox}
\small
A streaming service is \emph{free for the first 3 months}, then costs \emph{\$12 a month} for
every month after that. Let $C(m)$ be the \emph{total} you have paid after $m$ months.
\par\vspace{4pt}
Total paid after $2$ months: \blank{1.2cm} \quad after $5$ months: \blank{1.2cm} \quad
Could \emph{one} straight line describe the total for every month? \blank{1.0cm}
\par\vspace{5pt}
\textbf{Month 3 --- the last free month, or the first paid one?} Circle one: \quad
\textbf{free} \qquad \textbf{\$12} \qquad Why? \blank{5.2cm}
\par\vspace{4pt}
One function, two rules. Every question today is this picture: \textbf{which rule does this
month use --- and what does the graph do where the rule changes?}
\end{hookbox}

\small
% Displays inside table cells: tight skips, or every brace costs a centimetre. (\small resets
% the display skips, so this must follow it.) The work blocks use the same skips. Figures are
% centred with \centering inside the cell, not a center environment (its \topsep is ~1cm a figure).
\setlength{\abovedisplayskip}{4pt}\setlength{\belowdisplayskip}{4pt}%
\setlength{\abovedisplayshortskip}{2pt}\setlength{\belowdisplayshortskip}{2pt}%
\begin{guidednotes}
% ── ROW 1: one function, several rules — evaluate by choosing the piece ──────
\mainidea[One function,]{Several Rules} &
A \textbf{piecewise-defined function} uses different rules on different parts of its domain,
written with a brace. Each rule comes with the inputs it is \emph{for}.
\textbf{To evaluate: first pick the piece, then apply its rule.}
\\
& \[ C(m)=\begin{cases} 0, & 0\le m\le 3 \\ 12(m-3), & m>3 \end{cases} \]
{\centering
\begin{tikzpicture}[scale=0.52]
  \mgrid
  \draw[very thick, forest] (0,0)--(3,0);
  \draw[very thick, forest] (3,0)--(7,4);
  \fdot{forest}{3}{0}
  \node[above, font=\scriptsize\bfseries, forest] at (1.5,0.1) {$C(m)=0$};
  \node[above left, font=\scriptsize\bfseries, forest] at (5.6,2.6) {$C(m)=12(m-3)$};
\end{tikzpicture}\par\vspace{2pt}
The rule changes at $m=3$. That $m$-value is called the \labelbox{2.4cm}{}\par}
\\
& \textbf{Worked:} $C(5)$ --- is $5>3$? Yes, so use the second piece: $C(5)=12(5-3)=24$.
\par\vspace{2pt}
\textbf{Boundary rule:} an input on the boundary belongs to the piece whose inequality
\emph{includes} it --- the one with $\le$ or $\ge$, never the strict one. The story does not
decide it; the inequality does.
\\
& \notesprompt{Pick the piece, then apply its rule.}
\begin{probgrid}{|Y|Y|}
\pcell{1}{$C(2)$ --- is $2\le3$? Which piece, and what does it give?}{1.8cm}{} &
\pcell{2}{$C(3)$ --- the boundary. Which piece owns month~3, and how do you know?}{1.8cm}{} \\ \hline
\end{probgrid}
\\ \hline
% ── ROW 2: the V is two pieces ───────────────────────────────────────────────
\mainidea[The V is]{Two Pieces} &
\begin{minipage}[c]{0.58\linewidth}
Lesson~1.3's V is two lines pasted together at the vertex, and now we can \emph{write} it
that way. When $x\ge0$ the bars change nothing; when $x<0$ they flip the sign --- and the
check $|-3|=-(-3)=3$ shows the ``$-x$'' piece is not negative.
\[ |x|=\begin{cases} x, & x\ge0 \\ -x, & x<0 \end{cases} \]
\end{minipage}\hfill
\begin{minipage}[c]{0.38\linewidth}\centering
\begin{tikzpicture}[scale=0.36]
  \lgrid{-0.5}{6.5}{-0.5}{7.5}
  \draw[very thick, redacc] (0,7)--(3,1);
  \draw[very thick, forest] (3,1)--(6,7);
  \fdot{forest}{3}{1}
  \node[below, font=\tiny, charcoal] at (3,-0.05) {$3$};
  \node[left, font=\tiny, charcoal] at (-0.05,1) {$1$};
  \node[font=\scriptsize\bfseries] at (3,8.3) {$H(x)=2|x-3|+1$};
\end{tikzpicture}\par\vspace{2pt}
The vertex is at \labelbox{1.4cm}{}
\end{minipage}
\\
& \textbf{The trail from Lesson~1.3,} $H(x)=2|x-3|+1$, splits at its vertex $x=3$:
\par
\begin{minipage}[t]{0.47\linewidth}
For $x\ge3$, $|x-3|=x-3$:
\begin{work}
  H(x) &= 2(x-3) + 1 \\
       &= 2x - 5
\end{work}
\end{minipage}\hfill
\begin{minipage}[t]{0.47\linewidth}
For $x<3$, $|x-3|=-(x-3)=3-x$:
\begin{work}
  H(x) &= 2(3-x) + 1 \\
       &= 7 - 2x
\end{work}
\end{minipage}
\\
& \notesprompt{Use the two rules.}
\begin{probgrid}{|Y|Y|}
\pcell{3}{At the boundary, what does \emph{each} rule give for $H(3)$? So do the pieces
meet or jump there?}{1.8cm}{} &
\pcell{4}{Evaluate $H(1)$ two ways: pick the piece, then check against
$2|1-3|+1$.}{1.8cm}{} \\ \hline
\end{probgrid}
\\ \hline
% ── ROW 3: open or closed, meets or jumps ────────────────────────────────────
\mainidea[Open or closed,]{Meets or Jumps} &
A \textbf{closed} dot $\bullet$ \emph{includes} its endpoint ($\le,\ge$); an \textbf{open} dot
$\circ$ \emph{excludes} it ($<,>$). Same story, two different functions --- the \emph{total} paid
so far, and the \emph{price this month}:
\[ C(m)=\begin{cases} 0, & 0\le m\le3 \\ 12(m-3), & m>3 \end{cases} \qquad\qquad
   P(m)=\begin{cases} 0, & 0\le m\le3 \\ 12, & m>3 \end{cases} \]
\\
& {\centering
\begin{tikzpicture}[scale=0.46]
  \node[font=\small\bfseries] at (3.5,5.3) {Total paid $C(m)$};
  \mgrid
  \draw[very thick, forest] (0,0)--(3,0);
  \draw[very thick, forest] (3,0)--(7,4);
  \fdot{forest}{3}{0}
  \begin{scope}[xshift=9.6cm]
    \node[font=\small\bfseries] at (3.5,5.3) {Price this month $P(m)$};
    \mgrid
    \draw[very thick, redacc] (0,0)--(3,0);
    \draw[very thick, redacc] (3,1)--(7,1);
    \fdot{redacc}{3}{0}
    \hdot{redacc}{3}{1}
    \node[right, font=\tiny, charcoal] at (3.25,1.35) {$(3,12)$};
    \node[below right, font=\tiny, charcoal] at (3.15,-0.05) {$(3,0)$};
  \end{scope}
\end{tikzpicture}\par}
\vspace{2pt}
\textbf{One test:} both rules agree at the boundary, or they do not.
\\
& \notesprompt{Read the two graphs.}
\begin{probgrid}{|Y|Y|}
\pcell{5}{$C(m)$: at $m=3$ the first rule gives what, and just past $3$ the second? Meet or
jump? Constant on what interval, increasing on what?}{2.0cm}{} &
\pcell{6}{$P(m)$: the same questions --- and what is its range?}{2.0cm}{} \\ \hline
\end{probgrid}
\\ \hline
% ── ROW 4: who owns the boundary — THE CRUX, then rounding down ──────────────
\mainidea[Who owns the]{Boundary} &
\textbf{Exactly one piece owns a boundary input --- the piece whose inequality includes it.
The story does not decide it; the inequality does, and the closed dot shows which piece won.}
Above, the dot at $(3,0)$ is \labelbox{1.4cm}{} and the dot at $(3,12)$ is \labelbox{1.4cm}{}.
\\
& \begin{minipage}[c]{0.63\linewidth}
$P(m)$ is \emph{constant} on each piece, then jumps: a \textbf{step function}. The
\textbf{greatest-integer function} $\lfloor x\rfloor$ is the greatest integer \emph{less than or
equal to} $x$: it \textbf{rounds down}, and down means \emph{more negative}, not ``toward zero.''
On $[n,n+1)$ the output is the constant $n$ --- closed on the left, open on the right: the
boundary rule, one step at a time.
\end{minipage}\hfill
\begin{minipage}[c]{0.34\linewidth}\centering
\begin{tikzpicture}[scale=0.4]
  \node[font=\small\bfseries] at (1.5,4.1) {$y=\lfloor x\rfloor$};
  \lgrid{-1.5}{4.5}{-1.5}{3.6}
  \draw[very thick, forest] (-1,-1)--(0,-1); \fdot{forest}{-1}{-1} \hdot{forest}{0}{-1}
  \draw[very thick, forest] (0,0)--(1,0);    \fdot{forest}{0}{0}   \hdot{forest}{1}{0}
  \draw[very thick, forest] (1,1)--(2,1);    \fdot{forest}{1}{1}   \hdot{forest}{2}{1}
  \draw[very thick, forest] (2,2)--(3,2);    \fdot{forest}{2}{2}   \hdot{forest}{3}{2}
  \draw[very thick, forest] (3,3)--(4,3);    \fdot{forest}{3}{3}   \hdot{forest}{4}{3}
\end{tikzpicture}\par\vspace{2pt}
This graph is a \labelbox{2.4cm}{}
\end{minipage}
\\
& \notesprompt{Settle the vote.}
\begin{probgrid}{|Y|Y|}
\pcell{7}{A classmate reads $P(3)=12$ ``because the price is \$12 after the free months.''
What is $P(3)$, and which dot says so? What does the other dot mark?}{2.2cm}{} &
\pcell{8}{\textbf{In your own words:} what does a closed dot tell you, and what does an open dot
tell you?}{2.2cm}{} \\ \hline
\end{probgrid}
\\
& \textbf{9.} Round down --- fill the table. Then: does the step $[2,3)$ own $x=3$? \blank{1.6cm}
\par\vspace{2pt}
\renewcommand{\arraystretch}{1.2}
\begin{tabularx}{\linewidth}{@{} l Y Y Y Y @{}}
\rowcolor{sky}
\textbf{$x$} & $3.7$ & $5$ & $-0.4$ & $-2.5$ \\
\hline
\textbf{$\lfloor x\rfloor$} & \blank{0.9cm} & \blank{0.9cm} & \blank{0.9cm} & \blank{0.9cm} \\
\end{tabularx}
\\ \hline
% ── ROW 5: guided practice — WE DO, together, no story ───────────────────────
\mainidea[Guided practice]{Just the Rule} &
\begin{minipage}[c]{0.52\linewidth}
No story this time --- just the rule and its graph, pre-drawn.
\[ p(x)=\begin{cases} -2x, & x\le -1 \\ x+3, & x>-1 \end{cases} \]
\end{minipage}\hfill
\begin{minipage}[c]{0.44\linewidth}\centering
\begin{tikzpicture}[scale=0.3]
  \lgrid{-4.5}{4.5}{-1}{8.5}
  \draw[very thick, redacc] (-4,8)--(-1,2);
  \draw[very thick, forest] (-1,2)--(4,7);
  \fdot{forest}{-1}{2}
  \node[below, font=\tiny, charcoal] at (-1,-0.05) {$-1$};
  \node[left, font=\tiny, charcoal] at (-0.05,2) {$2$};
  \node[above left, font=\scriptsize\bfseries, redacc] at (-2.5,5) {$-2x$};
  \node[above left, font=\scriptsize\bfseries, forest] at (2.5,5.5) {$x+3$};
\end{tikzpicture}
\end{minipage}
\\
& \notesprompt{We work these together.}
\begin{probgrid}{|Y|Y|}
\pcell{10}{Pick the piece, then apply it: $p(-3)$, $p(0)$, $p(2)$ --- then $p(-1)$. Which piece
owns $x=-1$, and how do you know?}{1.8cm}{} &
\pcell{11}{Meets or jumps at $x=-1$? From the graph: decreasing where, increasing where, and
the minimum value?}{1.8cm}{} \\ \hline
\end{probgrid}
\\
& \begin{probgrid}*{|Y|Y|}
\pcell{12}{\textbf{The crux, on new ground.} Change the second piece to $x+4$. Now who owns
$x=-1$, what is $p(-1)$, meet or jump --- and which dot is closed? A classmate said ``both rules
gave $2$, so it does not matter who owns it'': what did they miss?}{2.6cm}{} &
\pcell{13}{\textbf{Round down:} $\lfloor 2.9\rfloor$, $\lfloor -1.2\rfloor$, $\lfloor 4\rfloor$.
Then $\lfloor x\rfloor=4$ for every $x$ in which interval --- and does that step own
$x=5$?}{2.6cm}{} \\ \hline
\end{probgrid}
\\ \hline
\end{guidednotes}

% ── THE NOTES END AT GUIDED PRACTICE ─────────────────────────────────────────
% There is deliberately no "you do" here: ../homework/ is the individual practice,
% started in class in the last ten minutes while the teacher circulates.

\end{document}
